\chapter{Symétrie, dégénérescence et échelles de moment cinétique}
\providecommand{\ket}[1]{\left\lvert#1\right\rangle}

\section{Symétries et niveaux d’énergie dégénérés}

Une symétrie est une opération qui ne modifie pas la description physique d’un système. Elle ne change ni le hamiltonien, ni l’équation de Schr\"odinger, ni les énergies possibles. Si un opérateur unitaire \(S\) représente l’opération, l’invariance signifie
\begin{equation}
S H S^\dagger=H,
\end{equation}
ou, de manière équivalente,
\begin{equation}
[S,H]=0.
\end{equation}

La translation en est un exemple. Déplacer un atome d’un endroit à un autre ne change pas son énergie. Cela vaut pour tous les états atomiques, et pas seulement pour l’état fondamental. La rotation en est un autre exemple, qui constituera le principal sujet de cette section.

L’échange de deux particules identiques est également une symétrie. Appeler deux électrons électron un et électron deux, ou Harry et Fred, ne les rend pas différents. L’échange des étiquettes ne peut pas modifier les propriétés de l’atome.

Un cristal possède une autre forme de symétrie de translation. Déplacer tout le cristal dans l’espace est une translation ordinaire. Toutefois, pour un cristal infini, une translation d’une période du réseau rétablit la même configuration physique. De manière équivalente, on peut renuméroter les sites du réseau en décalant d’une unité l’indice de chaque site. Les bords empêchent cette propriété d’être exacte pour un cristal fini, mais elle constitue une bonne approximation loin des extrémités d’un cristal très long.

L’homogénéité de l’espace---le fait que toute position ressemble à toute autre position---est la symétrie de translation. L’isotropie de l’espace---le fait que toute direction ressemble à toute autre direction---est la symétrie de rotation.

La dégénérescence est liée à la symétrie, mais ce n’est pas la même chose. Un niveau d’énergie est dégénéré lorsque plusieurs états possèdent exactement cette énergie. Pour un hamiltonien générique à spectre discret, une égalité exacte entre deux niveaux sans lien entre eux serait une coïncidence remarquable. Un accord à une partie sur \(10^{15}\) ne constitue toujours pas une égalité exacte. En général, les dégénérescences exactes signalent une symétrie.

La symétrie implique parfois une dégénérescence, mais pas toujours. Le premier exemple montre à la fois pourquoi c’est le cas et pourquoi une symétrie unique peut ne pas suffire.

\section{Une particule sur un cercle}

Considérons une bille contrainte à se déplacer sur un fil circulaire. On pourrait décrire sa position à l’aide des coordonnées cartésiennes \(x\) et \(y\), mais une seule coordonnée angulaire suffit :
\begin{equation}
\psi=\psi(\theta).
\end{equation}
La densité de probabilité de position sur le cercle est \(\psi^*(\theta)\psi(\theta)\).

Une rotation active dans le sens antihoraire d’un petit angle \(\epsilon\) transforme
\begin{equation}
\psi(\theta)\longrightarrow\psi(\theta-\epsilon).
\end{equation}
Par conséquent,
\begin{align}
\delta\psi(\theta)
&=\psi(\theta-\epsilon)-\psi(\theta)\\
&=-\epsilon\frac{d\psi}{d\theta}
+\mathcal O(\epsilon^2).
\end{align}

L’opérateur
\begin{equation}
L_z=-i\hbar\frac{d}{d\theta}
\end{equation}
est hermitien sur l’espace des fonctions d’onde périodiques. C’est l’analogue angulaire de
\begin{equation}
p_x=-i\hbar\frac{d}{dx}.
\end{equation}
Le facteur \(i\) est essentiel pour l’hermiticité. En termes de \(L_z\),
\begin{equation}
\delta\psi
=-\frac{i\epsilon}{\hbar}L_z\psi
+\mathcal O(\epsilon^2),
\end{equation}
l’opérateur de rotation infinitésimale est donc
\begin{equation}
R(\epsilon)
=1-\frac{i\epsilon}{\hbar}L_z
+\mathcal O(\epsilon^2).
\end{equation}
Cela identifie \(L_z\) comme le générateur des rotations autour de l’axe perpendiculaire au cercle.

\begin{classroomqa}
\lecturetimestamp{00:11:19}
\audiencequestion{L’un des facteurs faisant intervenir \(i\) ne devrait-il pas avoir le signe opposé ?}
\lecturerresponse{La vérification des signes donne la relation ci-dessus. Avec la convention \(\psi(\theta)\to\psi(\theta-\epsilon)\), la variation est \(-\epsilon\,d\psi/d\theta\), ce qui est égal à \(-i\epsilon L_z\psi/\hbar\).}
\end{classroomqa}

Cherchons maintenant les états propres du moment cinétique :
\begin{equation}
L_z\psi=M\psi.
\end{equation}
En substituant l’opérateur différentiel, on obtient
\begin{equation}
-i\hbar\frac{d\psi}{d\theta}=M\psi,
\end{equation}
dont les solutions sont de la forme
\begin{equation}
\psi(\theta)=C\exp\left(\frac{iM\theta}{\hbar}\right).
\end{equation}

Il existe une condition supplémentaire, absente du mouvement sur une droite infinie. Les points \(\theta=0\) et \(\theta=2\pi\) sont identiques ; une fonction d’onde orbitale univoque doit donc satisfaire
\begin{equation}
\psi(\theta+2\pi)=\psi(\theta).
\end{equation}
Pour la fonction propre ci-dessus, cela impose
\begin{equation}
\exp\left(\frac{2\pi iM}{\hbar}\right)=1.
\end{equation}
Par conséquent,
\begin{equation}
M=m\hbar,
\qquad
m\in\mathbb Z.
\end{equation}
Les valeurs possibles du moment cinétique orbital sur le cercle sont donc
\begin{equation}
\ldots,-2\hbar,-\hbar,0,\hbar,2\hbar,\ldots.
\end{equation}

Un état de moment cinétique \(m\hbar\) doit-il avoir la même énergie qu’un état de moment cinétique \(-m\hbar\) ? Les valeurs positive et négative de \(m\) décrivent des sens de circulation opposés. La question est de savoir si la symétrie de rotation, à elle seule, impose que ces deux mouvements aient la même énergie.

\begin{classroomqa}
\lecturetimestamp{00:22:51}
\audiencequestion{La convention \(\psi(\theta-\epsilon)\) signifie-t-elle que la particule se déplace dans le sens horaire ?}
\lecturerresponse{Non. Cette convention fixe uniquement le sens appelé rotation positive. Elle définit le signe de la position angulaire et celui du moment cinétique, mais elle ne détermine pas le sens réel du déplacement de la particule.}
\end{classroomqa}

Un champ magnétique perpendiculaire au plan du cercle fournit un contre-exemple. Il peut décaler les énergies des moments cinétiques positifs et négatifs en sens opposés, et ainsi lever la dégénérescence de la paire. Le signe du décalage dépend des conventions choisies pour la charge et l’orientation, mais c’est la levée de dégénérescence qui importe.

Un tel champ préserve toujours les rotations autour de l’axe perpendiculaire. Il montre donc que la symétrie de rotation, à elle seule, n’impose pas
\begin{equation}
E(m)=E(-m).
\end{equation}
En l’absence de cette levée de dégénérescence, \(m=1\) et \(m=-1\) peuvent former une paire dégénérée, tout comme \(m=2\) et \(m=-2\). L’état \(m=0\) n’a pas de partenaire distinct de ce type.

\section{Symétrie de réflexion et paire \texorpdfstring{\(m,-m\)}{m, -m}}

Ajoutons une symétrie discrète de réflexion. Désignons par \(\mathcal M\) la réflexion par rapport à un diamètre du cercle :
\begin{equation}
(\mathcal M\psi)(\theta)=\psi(-\theta).
\end{equation}
Elle inverse le sens de circulation :
\begin{equation}
\mathcal M\psi_m=\psi_{-m}.
\end{equation}

Si la réflexion est une symétrie, alors
\begin{equation}
[\mathcal M,H]=0.
\end{equation}
Supposons que
\begin{equation}
H\ket{m}=E_m\ket{m}.
\end{equation}
Appliquons \(\mathcal M\) :
\begin{align}
H\mathcal M\ket{m}
&=\mathcal M H\ket{m}\\
&=E_m\mathcal M\ket{m}.
\end{align}
Puisque \(\mathcal M\ket{m}\) est l’état de moment cinétique \(-m\hbar\),
\begin{equation}
E_{-m}=E_m.
\end{equation}
La symétrie de rotation conserve le moment cinétique ; la symétrie de réflexion relie les deux valeurs opposées. Ensemble, elles imposent que la paire soit dégénérée.

\begin{classroomqa}
\lecturetimestamp{00:29:54}
\audiencequestion{Puisque les modes \(m\) et \(-m\) ont la même densité de probabilité, cela suffit-il à montrer que leurs énergies sont égales ?}
\lecturerresponse{Non. Les fonctions d’onde \(e^{im\theta}\) et \(e^{-im\theta}\) sont différentes, même si leurs modules au carré sont égaux. L’équation aux valeurs propres du moment cinétique porte sur la fonction d’onde, pas seulement sur \(\psi^*\psi\). La mécanique quantique requiert les amplitudes, et pas uniquement les probabilités.}
\end{classroomqa}

Le champ magnétique préserve les rotations autour de son axe, mais brise la réflexion utilisée ci-dessus.

\begin{classroomqa}
\lecturetimestamp{00:33:05}
\audiencequestion{La réflexion inverse-t-elle le champ magnétique ?}
\lecturerresponse{Oui. Un champ magnétique est un vecteur axial. La réflexion d’un champ dirigé vers l’intérieur de la page donne un champ dirigé vers l’extérieur. De manière équivalente, une boucle de courant et son image miroir portent des courants de sens opposés et produisent des champs de directions opposées.}
\end{classroomqa}

La non-commutativité de la rotation et de la réflexion se voit directement. Sur une fonction propre du moment cinétique,
\begin{align}
\mathcal M L_z\psi_m
&=m\hbar\,\psi_{-m},\\
L_z\mathcal M\psi_m
&=-m\hbar\,\psi_{-m}.
\end{align}
Ainsi,
\begin{equation}
[\mathcal M,L_z]\psi_m
=2m\hbar\,\psi_{-m},
\end{equation}
ce qui est non nul pour \(m\neq0\). Effectuer une rotation puis une réflexion n’est pas la même opération qu’effectuer une réflexion puis une rotation.

Lorsque les deux transformations sont des symétries, leurs opérateurs commutent tous deux avec \(H\), mais ils ne commutent pas entre eux. C’est cette structure qui est à l’origine de la dégénérescence.

\begin{classroomqa}
\lecturetimestamp{00:40:28}
\audiencequestion{Toute paire non commutative d’opérateurs de Heisenberg, telle que \(x\) et \(p\), engendre-t-elle des dégénérescences ?}
\lecturerresponse{En général, non. Les opérateurs doivent aussi être des symétries : ils doivent commuter avec le hamiltonien. En général, la position et la quantité de mouvement ne sont pas toutes deux conservées. Si elles commutaient toutes deux avec \(H\), leur non-commutativité impliquerait une dégénérescence.}
\end{classroomqa}

\section{Algèbres de commutateurs des symétries}

Supposons que \(A\) et \(B\) soient les générateurs hermitiens de deux symétries continues :
\begin{equation}
[A,H]=0,
\qquad
[B,H]=0.
\end{equation}
Rien dans ces équations n’exige que \(A\) et \(B\) commutent entre eux.

Le commutateur \([A,B]\) est antihermitien ; définissons donc un opérateur hermitien
\begin{equation}
C=i[A,B].
\end{equation}
Les facteurs numériques et le signe global sont sans importance ici. L’essentiel est que \(C\) commute également avec le hamiltonien. En effet,
\begin{equation}
[AB,H]
=A[B,H]+[A,H]B=0,
\end{equation}
et, de même,
\begin{equation}
[BA,H]=0.
\end{equation}
Par conséquent,
\begin{equation}
[C,H]
=i\bigl([AB,H]-[BA,H]\bigr)=0.
\end{equation}

Le nouveau générateur \(C\) peut être indépendant de \(A\) et \(B\), ou n’être qu’une combinaison linéaire de générateurs déjà connus. S’il est nouveau, ajoutons-le à la liste et formons d’autres commutateurs. Ce processus peut se poursuivre indéfiniment, mais il s’arrête souvent après un nombre fini de générateurs : tout commutateur supplémentaire est une combinaison linéaire de ceux déjà présents. Ce système fermé constitue une algèbre de commutateurs, ou algèbre de Lie.

Les transformations de symétrie finies forment un groupe. Leurs générateurs infinitésimaux forment l’algèbre de Lie associée.

Les dégénérescences sont importantes parce qu’elles organisent les spectres. Si deux états ont la même énergie, l’un ne peut pas se désexciter vers l’autre en émettant un photon transportant la différence de leurs énergies. En spectroscopie atomique, la disposition et la dégénérescence des niveaux déterminent les raies spectrales possibles. Le même raisonnement s’applique aux hadrons excités et à leurs transitions.

\begin{classroomqa}
\lecturetimestamp{00:49:07}
\audiencequestion{Ces commutateurs sont-ils eux-mêmes des éléments du groupe de symétrie ?}
\lecturerresponse{Les opérateurs dont il est question sont des générateurs. Les générateurs et leurs commutateurs forment une algèbre de commutateurs, ou algèbre de Lie. Les opérations de symétrie finies construites à partir des générateurs sont les éléments du groupe.}
\end{classroomqa}

Si tous les générateurs de symétrie commutent entre eux, la symétrie est dite abélienne. Si certains de leurs commutateurs sont non nuls, elle est non abélienne.

\begin{classroomqa}
\lecturetimestamp{00:50:25}
\audiencequestion{Quelle est la définition fondamentale d’un générateur ?}
\lecturerresponse{Considérons une transformation de symétrie unitaire très proche de l’identité et écrivons-la sous la forme \(1-i\epsilon G/\hbar+\cdots\). L’opérateur hermitien \(G\) en est le générateur. La répétition de la transformation infinitésimale correspondante construit une transformation finie.}
\end{classroomqa}

\begin{classroomqa}
\lecturetimestamp{00:52:08}
\audiencequestion{Que désigne l’algèbre des générateurs de symétrie ?}
\lecturerresponse{Il s’agit d’un ensemble de générateurs que l’on peut additionner, soustraire, multiplier par des nombres ordinaires et commuter. Dans cette algèbre, le commutateur joue le rôle de la multiplication ; la fermeture signifie que ces opérations ne produisent pas de générateurs extérieurs à l’ensemble.}
\end{classroomqa}

\section{Rotations autour d’axes différents}

Les rotations fournissent l’exemple le plus important de symétrie non abélienne. Une rotation autour de l’axe \(x\) suivie d’une rotation autour de l’axe \(y\) n’est généralement pas équivalente à ces rotations effectuées dans l’ordre inverse. La différence entre les deux est elle-même liée à une rotation autour du troisième axe.

\begin{classroomqa}
\lecturetimestamp{00:54:22}
\audiencequestion{Les rotations infinitésimales commutent-elles, même si les rotations finies ne commutent pas ?}
\lecturerresponse{Elles diffèrent au second ordre par rapport aux deux petits angles de rotation. Cela suffit : une fois ces petits paramètres mis en facteur, les générateurs eux-mêmes ont un commutateur non nul.}
\end{classroomqa}

Avant de calculer les commutateurs quantiques, considérons l’image classique correspondante. Prenons une orbite planétaire circulaire dont le moment cinétique est perpendiculaire au plan. Faire tourner toute l’orbite autour de l’axe \(y\) donne une autre orbite admissible, d’énergie exactement identique. Le vecteur moment cinétique a toutefois maintenant une autre composante \(x\).

L’invariance par rotation produit donc des configurations de même énergie, mais dont les composantes du moment cinétique diffèrent. En mécanique quantique, cela se manifeste par la non-commutativité des générateurs de rotations autour d’axes différents.

\section{Le moment cinétique en coordonnées cartésiennes}

Considérons d’abord une particule qui se déplace dans le plan \(xy\), sans la contraindre à rester sur un cercle. Sous une petite rotation antihoraire,
\begin{equation}
\delta x=-\epsilon y,
\qquad
\delta y=\epsilon x.
\end{equation}
Pour l’état tourné activement, la fonction d’onde transformée est évaluée aux coordonnées tournées de manière inverse :
\begin{equation}
\psi_\epsilon(x,y)
=\psi(x+\epsilon y,y-\epsilon x).
\end{equation}
Par conséquent,
\begin{align}
\delta\psi
&=\epsilon\left(
y\frac{\partial}{\partial x}
-x\frac{\partial}{\partial y}
\right)\psi
+\mathcal O(\epsilon^2)\\
&=-\frac{i\epsilon}{\hbar}
\left(xp_y-yp_x\right)\psi
+\mathcal O(\epsilon^2).
\end{align}
Le générateur des rotations autour de l’axe \(z\) est donc
\begin{equation}
L_z=xp_y-yp_x.
\end{equation}
Il s’agit de la composante \(z\) du vecteur classique familier
\begin{equation}
\mathbf L=\mathbf r\times\mathbf p.
\end{equation}
Ce n’est pas seulement le moment cinétique mécanique ; en mécanique quantique, c’est l’opérateur qui engendre les rotations de la fonction d’onde.\footnote{La convention de rotation active est choisie de manière à être cohérente avec la transformation précédente \(\psi(\theta)\mapsto\psi(\theta-\epsilon)\) et avec l’identification finale \(L_z=xp_y-yp_x\), après l’incertitude intermédiaire sur le signe dans la démonstration orale.}

En appliquant cycliquement \(x\to y\to z\to x\), on obtient les autres composantes.\footnote{Les formules de \(L_x\) et \(L_y\) sont reconstituées en appliquant à \(L_z=xp_y-yp_x\) la substitution cyclique \(x\to y\to z\to x\), explicitement prescrite dans le cours.}
\begin{align}
L_x&=yp_z-zp_y,\\
L_y&=zp_x-xp_z,\\
L_z&=xp_y-yp_x.
\end{align}

Pour un système invariant par rotation, ces opérateurs commutent avec le hamiltonien :
\begin{equation}
[H,L_x]=[H,L_y]=[H,L_z]=0.
\end{equation}
Une particule dans un champ de force central en est l’exemple classique. Ces équations expriment la conservation quantique du moment cinétique.

\section{L’algèbre de Lie du moment cinétique}

Les relations de commutation canoniques sont
\begin{align}
[x_i,x_j]&=0,\\
[p_i,p_j]&=0,\\
[x_i,p_j]&=i\hbar\,\delta_{ij}.
\end{align}
Les coordonnées commutent avec les composantes de la quantité de mouvement selon d’autres directions. Ainsi, \(x\) commute avec \(p_y\), par exemple, tandis que \(x\) ne commute pas avec \(p_x\).

Calculons maintenant
\begin{equation}
[L_x,L_y]
=
[yp_z-zp_y,\;zp_x-xp_z].
\end{equation}
La plupart des termes s’annulent. Les deux contributions non nulles sont
\begin{align}
[yp_z,zp_x]&=-i\hbar\,yp_x,\\
[zp_y,xp_z]&=i\hbar\,xp_y.
\end{align}
Par conséquent,
\begin{align}
[L_x,L_y]
&=i\hbar\left(xp_y-yp_x\right)\\
&=i\hbar L_z.
\end{align}
La permutation cyclique des indices des coordonnées donne
\begin{align}
[L_x,L_y]&=i\hbar L_z,\\
[L_y,L_z]&=i\hbar L_x,\\
[L_z,L_x]&=i\hbar L_y.
\end{align}

\begin{classroomqa}
\lecturetimestamp{01:14:50}
\audiencequestion{Quand le commutateur donne \(L_z\), faut-il encore montrer que \(L_z\) n’est pas une combinaison linéaire de \(L_x\) et \(L_y\) ?}
\lecturerresponse{Oui. Il faut établir cette indépendance au lieu de la supposer. À ce stade, on ne cherche pas à en donner une démonstration algébrique directe.}
\end{classroomqa}

Les commutateurs affichés se referment sur l’espace engendré par \(L_x\), \(L_y\) et \(L_z\) : une fois ces trois opérateurs inclus, les commutateurs suivants ne produisent aucun opérateur extérieur à cet espace. C’est l’algèbre de Lie du moment cinétique. Avec
\begin{equation}
[H,L_i]=0,
\qquad i=x,y,z,
\end{equation}
ces relations livrent des informations substantielles sur le spectre d’énergie sans qu’il soit nécessaire de résoudre directement l’équation de Schr\"odinger angulaire.

Une stratégie algébrique semblable est utilisée pour l’oscillateur harmonique quantique, dont le spectre est engendré par les opérateurs de création et d’annihilation. Ici, les opérateurs analogues modifieront la composante \(z\) du moment cinétique plutôt que l’énergie.

Choisir l’axe \(z\) n’est qu’un choix de représentation. Rien ne rend \(z\) physiquement préférable à \(x\) ou à \(y\), mais choisir une composante fournit une base commode. Écrivons
\begin{equation}
L_z\ket{m}=m\hbar\ket{m}.
\end{equation}
Il s’agit de déterminer les valeurs que peut prendre \(m\) et comment les états ayant différentes valeurs de \(m\) sont reliés.

\begin{classroomqa}
\lecturetimestamp{01:24:02}
\audiencequestion{Quand vous comparez cette construction à celle de l’oscillateur harmonique, parlez-vous de l’oscillateur harmonique quantique ? Il n’a pas été abordé au cours précédent.}
\lecturerresponse{Oui, la comparaison porte sur l’oscillateur quantique. Le raisonnement sur le moment cinétique présenté ici se suffit à lui-même et ne nécessite aucune connaissance de l’oscillateur. L’analogie peut ensuite être comprise dans l’autre sens : l’oscillateur fait appel à une méthode algébrique similaire.}
\end{classroomqa}

\section{Augmentation et diminution du moment cinétique}

Définissons
\begin{equation}
L_+=L_x+iL_y,
\qquad
L_-=L_x-iL_y.
\end{equation}
Les relations de commutation du moment cinétique impliquent
\begin{equation}
[L_z,L_+]=\hbar L_+,
\qquad
[L_z,L_-]=-\hbar L_-.
\end{equation}

Appliquons la première relation à un état propre de \(L_z\) :
\begin{align}
L_zL_+\ket{m}
&=\left(L_+L_z+[L_z,L_+]\right)\ket{m}\\
&=(m+1)\hbar\,L_+\ket{m}.
\end{align}
À moins que \(L_+\ket{m}=0\), c’est un autre état propre de \(L_z\), dont la valeur propre est maintenant \((m+1)\hbar\). De même,
\begin{equation}
L_zL_-\ket{m}
=(m-1)\hbar\,L_-\ket{m}.
\end{equation}
Ainsi, \(L_+\) augmente \(m\) d’une unité et \(L_-\) le diminue d’une unité.

À partir d’une valeur propre quelconque, l’algèbre engendre une suite dont les termes successifs diffèrent d’un entier. La suite peut se prolonger indéfiniment ou s’arrêter parce que
\begin{equation}
L_+\ket{m_{\mathrm{top}}}=0
\end{equation}
à son extrémité supérieure, et que
\begin{equation}
L_-\ket{m_{\mathrm{bottom}}}=0
\end{equation}
à son extrémité inférieure.

\begin{classroomqa}
\lecturetimestamp{01:35:11}
\audiencequestion{L’argument du cercle n’imposait-il pas déjà que le moment cinétique soit entier ?}
\lecturerresponse{Cet argument établit que les valeurs du moment cinétique orbital représenté par des fonctions d’onde de position univoques sont entières. L’algèbre abstraite du moment cinétique donne un résultat moins contraignant : l’échelle peut contenir des valeurs entières ou demi-entières.}
\end{classroomqa}

Si l’échelle s’arrête, l’invariance par rotation relie ses deux extrémités. Une rotation de \(180^\circ\) autour d’un axe perpendiculaire à \(z\) transforme
\begin{equation}
L_z\longrightarrow-L_z.
\end{equation}
Le spectre doit donc être symétrique par rapport à zéro :
\begin{equation}
m_{\mathrm{bottom}}=-m_{\mathrm{top}}.
\end{equation}
Comme les valeurs voisines diffèrent d’une unité, deux possibilités se présentent. L’échelle contient soit les entiers,
\begin{equation}
\ldots,-2,-1,0,1,2,\ldots,
\end{equation}
soit les demi-entiers,
\begin{equation}
\ldots,-\frac{5}{2},-\frac{3}{2},-\frac{1}{2},
\frac{1}{2},\frac{3}{2},\frac{5}{2},\ldots,
\end{equation}
avec des extrémités auxquelles l’opérateur d’augmentation ou de diminution donne zéro.

L’argument fondé sur l’univocité de la fonction d’onde sur le cercle sélectionne, pour le moment cinétique orbital, la possibilité des valeurs entières :
\begin{equation}
\mathbf L=\mathbf r\times\mathbf p.
\end{equation}
La possibilité des valeurs demi-entières est également physique ; elle correspond au moment cinétique de spin, qui sera étudié séparément.

Enfin, les états de l’échelle ont la même énergie lorsque le hamiltonien est invariant par rotation. Supposons que
\begin{equation}
H\ket{m}=E\ket{m}.
\end{equation}
Puisque chaque composante \(L_i\) commute avec \(H\), \(L_+\) et \(L_-\) commutent également avec lui. On a donc
\begin{align}
H L_+\ket{m}
&=L_+H\ket{m}\\
&=E L_+\ket{m}.
\end{align}
Lorsque \(L_+\ket{m}\neq0\), l’état obtenu après augmentation a la même énergie que l’état initial. Le même calcul donne
\begin{equation}
H L_-\ket{m}=E L_-\ket{m}.
\end{equation}
L’augmentation et la diminution engendrent donc une famille d’états de valeurs propres \(L_z\) différentes, mais ayant exactement la même énergie. C’est le lien entre symétries non commutatives et dégénérescence, développé tout au long du cours.

\begin{classroomqa}
\lecturetimestamp{01:44:13}
\audiencequestion{Ces états de l’échelle sont-ils associés au choix de différents axes de rotation ?}
\lecturerresponse{Ils sont reliés par des rotations. Les vecteurs propres de \(L_x\) ou de \(L_y\) sont des combinaisons linéaires des vecteurs propres de \(L_z\) choisis. La rotation d’un état modifie ses composantes de moment cinétique tout en laissant son énergie inchangée.}
\end{classroomqa}

Le choix de \(L_+\) et \(L_-\) est commode, car ces combinaisons font apparaître directement l’échelle. Cette astuce revient souvent en mécanique quantique : lorsqu’une algèbre de commutateurs se referme, les opérateurs d’augmentation et de diminution organisent souvent le spectre des valeurs propres.

\begin{classroomqa}
\lecturetimestamp{01:46:33}
\audiencequestion{Que se passe-t-il lorsque \(L_+\) agit sur l’état maximal, et comment définit-on le sommet de l’échelle ?}
\lecturerresponse{Il donne zéro. C’est la définition du sommet d’une échelle qui s’arrête. De même, \(L_-\) donne zéro au bas de l’échelle. Une échelle abstraite n’est pas nécessairement bornée. Pour les états de moment cinétique atomique étudiés plus loin, on rencontre généralement des échelles qui s’arrêtent.}
\end{classroomqa}

L’étape suivante consiste à approfondir cette structure du moment cinétique et à l’appliquer aux atomes, y compris à l’hydrogène. Les multiplets de rotation établis ici rendent compte d’une partie du schéma de dégénérescence ; d’autres éléments de structure restent à développer.
